% ---------------------------------------------------------------------------
% Front matter: preface, how to read this book, the two pilgrims
% ---------------------------------------------------------------------------
\chapter{Preface: A Jet Is a Journey, Not a Destination}
\chaptermark{Preface}
\label{ch:preface}

In CMS an electron, a photon or a muon is a \emph{destination}: an object that
is essentially there at the end of reconstruction and needs only identification
and calibration. A jet is a \emph{journey}. It does not exist at any single
point; it is the accumulated history of a coloured parton from the hard
scattering to the calibrated, tagged, smeared four-vector an analyst finally
puts into a histogram. Every stage of that journey either adds energy to the
jet, removes energy from it, or reshuffles how the energy is distributed.

This book follows one and the same jet through every stage, shows the
mathematics of each stage, measures how much energy that stage added or
removed, and points to the official paper for further reading. The title is
meant literally: \emph{Journey Is The Jet}, and the two jets we follow are the
$b$-jet (bottom) and the top-jet (top).

\todo{Expand the preface: who this book is for (new Master, PhD and postdoc
members of the CMS JME group), what is assumed, what is not.}

\section*{How to read this book}

Chapters are stages of a pilgrimage told through one metaphor: a mother quark,
her children from the parton shower, the crowd on the road (pileup), borders and
checkpoints (detector and trigger), the census (Particle Flow), the police
(pileup mitigation), the inn (clustering), the tax office (jet energy
corrections), the wobbly ruler (resolution) and the registry office (tagging).
The metaphor is a memory aid; the physics is always stated explicitly.

\paragraph{The voice of this book.} Think of the journey as a guided tour.
The author is the guide and you are the passenger: at every stop the guide
explains what is happening outside the window, hands you the formula that
describes it, and asks you to press a button and see the number come out
yourself. The tone is deliberately that of a teacher standing beside you
rather than a reference manual in front of you, the book says so
and does it with you. The code that produces every number is not printed in
the book. It lives in the public repository
\url{https://github.com/ravindkv/JitJet}, and each chapter gives the command
to run and the output to expect, so that the pages can be spent on the
mathematics instead.

Every chapter has the same internal structure:
\begin{enumerate}[nosep]
  \item \textbf{The stage in the journey}: the metaphor and what physically happens.
  \item \textbf{The mathematics}: the minimal, correct formulae.
  \item \textbf{Our jet at this stage}: the numbers from the energy ledger, plots and event displays.
  \item \textbf{How CMS does it in practice}: configurations, \CMSSW modules, JME tools, NanoAOD branches.
  \item \textbf{Pitfalls and FAQs} for newcomers.
  \item \textbf{Further reading}: official CMS papers, detector-performance notes, theory papers.
\end{enumerate}

\section*{The two pilgrims}

\begin{description}
  \item[Event A, the $b$-jet.] \code{generate p p > b b\~{}} at $\sqrt{s}=13.6\TeV$ in
    \MG with a fixed random seed, showered and hadronised with \PYTHIA (CP5 tune),
    premixed with Run~3 pileup, simulated with \GEANT, passed through Level-1 and
    High-Level Trigger emulation, reconstructed with Particle Flow, \PUPPI,
    \antikt $R=0.4$ clustering, jet energy corrections and resolution smearing,
    and finally $b$-tagged.
  \item[Event B, the top-jet.] \code{generate p p > t t\~{}} with a boosted
    hadronically decaying top quark ($\pt>400\GeV$), the same chain,
    \antikt $R=0.8$ clustering, soft drop, $N$-subjettiness and a
    machine-learned top tagger.
\end{description}

\section*{The energy ledger}

For both jets we keep an \emph{energy ledger}: one row per stage with the jet
\pt, mass, constituent multiplicity and the change with respect to the previous
stage. The ledger is the quantitative red thread of the book and is closed in
\cref{ch:energy_ledger}.

\section*{Reproducibility}

All scripts, generator configurations, logs and event records used in the
book are in the \code{example/} directory of
\url{https://github.com/ravindkv/JitJet}, one subdirectory per stage; the
animated views of the calculations are collected in the JitJet playlist at
\url{https://www.youtube.com/playlist?list=PLFtvDH9g5Km0}.
\todo{Software releases, generator versions, random seeds and the container
image.}
